
\chapter{System Design and Implementation}\label{ch:system}

\section{Architecture Overview}

The Flower FRL Benchmark is structured as a Python package (\texttt{frl\_benchmark})
that integrates with the Flower federated learning framework~\cite{beutel2020flower}.
The system consists of three major components:

\begin{enumerate}
    \item A \textbf{server application} (\texttt{server\_app.py}) that runs the
          Flower \texttt{ServerApp}, manages the shared policy, and dispatches
          gradient aggregation to a pluggable strategy.
    \item A \textbf{client application} (\texttt{client\_app.py}) that runs the
          Flower \texttt{ClientApp} on each simulated worker, collecting trajectories
          and computing local policy gradient estimates.
    \item A \textbf{web dashboard} (\texttt{dashboard/}) providing a browser-based
          interface for experiment configuration and real-time monitoring.
\end{enumerate}

Figure~\ref{fig:architecture} illustrates the high-level interaction between these
components. In simulation mode, Flower spawns all workers as lightweight virtual
nodes within a single process, communicating via in-process message passing. In
deployment mode, workers run on separate machines and communicate over gRPC.

\begin{figure}[htbp]
    \centering
    \begin{tikzpicture}[
        node distance=1.6cm and 2.2cm,
        box/.style={draw, rounded corners, minimum width=2.8cm, minimum height=1cm,
                    text centered, font=\small},
        arr/.style={->, >=Stealth, thick}
    ]
        \node[box, fill=blue!10] (server) {Server (\texttt{FRLStrategy})};
        \node[box, fill=green!10, below left=of server]  (w1) {Worker 1};
        \node[box, fill=green!10, below=of server]       (w2) {Worker 2};
        \node[box, fill=green!10, below right=of server] (wk) {Worker $K$};
        \node[box, fill=orange!10, right=of server]      (dash) {Dashboard};

        \draw[arr] (server) -- node[left,  font=\scriptsize]{params} (w1);
        \draw[arr] (server) -- node[right, font=\scriptsize]{params} (w2);
        \draw[arr] (server) -- node[right, font=\scriptsize]{params} (wk);
        \draw[arr] (w1) -- node[left,  font=\scriptsize]{grads} (server);
        \draw[arr] (w2) -- node[right, font=\scriptsize]{grads} (server);
        \draw[arr] (wk) -- node[right, font=\scriptsize]{grads} (server);
        \draw[arr, dashed] (server) -- node[above, font=\scriptsize]{metrics} (dash);
    \end{tikzpicture}
    \caption{High-level architecture of the Flower FRL Benchmark. The server broadcasts
    policy parameters to $K$ workers, collects gradient estimates, aggregates them via
    the selected strategy, and pushes metrics to the web dashboard.}
    \label{fig:architecture}
\end{figure}

\section{Flower Integration}

\subsection{Server Strategy}

The server implements Flower's \texttt{Strategy} abstract class through
\texttt{FRLStrategy}. The key methods are:

\begin{itemize}
    \item \texttt{configure\_fit}: Broadcasts current policy parameters to all
          workers and specifies the batch size for this round.
    \item \texttt{aggregate\_fit}: Collects gradient submissions, removes workers
          that returned zero examples (skipped rounds), reconstructs gradient
          tensors, and delegates to the selected aggregation strategy.
    \item \texttt{evaluate}: Evaluates the current policy on 10 episodes using
          the server's own environment instance and pushes the result to the dashboard.
\end{itemize}

The server uses a separate \texttt{gymnasium} environment instance for SCSG
variance-reduction steps (required by FedPG-BR and SVRPG), keeping the
trajectory sampling for gradient computation entirely local to the server.

\subsection{Client Application}

Each worker is represented by a \texttt{FRLClient}, which wraps a \texttt{Worker}
object. In each round, the client:
\begin{enumerate}
    \item Receives the current global policy parameters from the server.
    \item Overwrites its local policy with the received parameters.
    \item Samples $B$ trajectories from its local environment, computing the
          REINFORCE gradient estimate.
    \item Returns the gradient as a parameter array together with diagnostic
          metrics.
\end{enumerate}

Byzantine workers override the gradient computation to substitute a malicious
gradient, as described in \Cref{sec:attacks}.

\subsection{Simulation and Deployment}

In simulation mode the entire experiment runs as a single \texttt{flower-simulation}
process. All workers are spawned as virtual nodes; communication happens in memory.
This mode is the primary interface exposed through the web dashboard.

For real distributed deployment, \texttt{deploy\_server.py} and
\texttt{deploy\_client.py} provide entry points that connect to a Flower
superlink. This allows the benchmark to be used across multiple physical machines,
reflecting the intended federated setting.

\section{The Strategy Plugin System}

\subsection{Design}

The central design goal of the benchmark is extensibility: a researcher should be
able to contribute a new aggregation strategy without modifying any existing code.
This is achieved through a \emph{registry pattern} combined with a Python
decorator.

Each strategy is a class that inherits from \texttt{AggregationStrategy} and
implements two methods:

\begin{itemize}
    \item \texttt{aggregate(gradients, batch\_size, **kwargs)}: Takes the list of
          gradient tensors from all active workers and returns an aggregated gradient
          $\mu_t$ and the list of accepted worker indices.
    \item \texttt{server\_update(policy, optimizer, theta\_t\_0, mu\_t, config,
          env)}: Applies the aggregated gradient to the policy, optionally running
          multiple variance-reduction steps.
\end{itemize}

Registering a strategy requires only applying the \texttt{@register\_strategy}
decorator:

\begin{verbatim}
from frl_benchmark.strategies.base import (
    AggregationStrategy, register_strategy, apply_gradient
)

@register_strategy("my-method")
class MyStrategy(AggregationStrategy):
    def aggregate(self, gradients, batch_size, **kwargs):
        mu_t = torch.mean(torch.stack(gradients), dim=0)
        return mu_t, list(range(len(gradients)))

    def server_update(self, policy, optimizer,
                      theta_t_0, mu_t, config, env=None, env_name=""):
        apply_gradient(policy, optimizer, mu_t)
        return 1
\end{verbatim}

Once registered, the strategy is immediately selectable from the dashboard's
method dropdown without any further changes.

\subsection{Auto-Discovery}

Strategies are discovered at import time via \texttt{pkgutil.iter\_modules}, which
scans the \texttt{frl\_benchmark/strategies/} directory for Python modules and
imports them. This means a researcher can place a new \texttt{.py} file in that
directory and it will be found automatically on the next run.

\subsection{Bundled Strategies}

Three strategies are provided out of the box:

\begin{description}
    \item[\texttt{fedpg-br}] Full FedPG-BR: Byzantine filtering followed by SCSG
         variance reduction. The Byzantine filter is parameterised by $\sigma$ and
         $\delta$; SCSG uses a geometric distribution over the number of inner steps.
    \item[\texttt{svrpg}] SVRPG~\cite{papini2018svrpg}: SCSG variance reduction
         without Byzantine filtering. Accepts all workers unconditionally.
    \item[\texttt{gomdp}] GOMDP: Simple gradient averaging with a single gradient
         step. No filtering, no variance reduction. Serves as the simplest baseline.
\end{description}

An \texttt{example\_strategy.py} file demonstrates a complete Trimmed Mean
implementation, and \texttt{my\_strategy.py} provides a template with placeholder
methods for users to fill in.

\section{Environments and Configuration}

\subsection{Supported Environments}

The benchmark supports the following environments from Gymnasium:

\begin{table}[htbp]
\centering
\begin{tabular}{llll}
\hline
\textbf{Environment} & \textbf{Actions} & \textbf{State dim} & \textbf{Notes} \\
\hline
CartPole-v1       & Discrete   & 4  & Simplest testbed; bounded episodes \\
Acrobot-v1        & Discrete   & 6  & Harder; sparse reward \\
MountainCar-v0    & Discrete   & 2  & Very sparse reward \\
LunarLander-v3    & Discrete   & 8  & Variable episode length \\
HalfCheetah-v5    & Continuous & 17 & Requires MuJoCo \\
\hline
\end{tabular}
\caption{Supported Gymnasium environments and their properties.}
\label{tab:envs}
\end{table}

Each environment has a default configuration stored in \texttt{config.py},
including network architecture (hidden units, activation), learning rate,
batch size, and filter parameters $\sigma$ and $\delta$. All parameters
can be overridden from the dashboard's advanced settings panel.

\subsection{Policy Networks}

Discrete-action environments use a softmax MLP policy; continuous-action
environments use a Gaussian MLP policy with a learned diagonal covariance.
All policies are implemented in \texttt{policy.py} and support \texttt{get\_flat\_params}
and \texttt{set\_flat\_params} for efficient parameter serialisation over the
Flower communication layer.

\section{Byzantine Attack Simulation}\label{sec:attacks}

The benchmark simulates seven attack types to enable evaluation of Byzantine
robustness:

\begin{enumerate}
    \item \textbf{Random Noise (RN)}: Replaces the gradient with a random vector
          drawn from a standard normal distribution.
    \item \textbf{Random Action (RA)}: The worker takes uniformly random actions
          during trajectory collection, simulating sensor or hardware failure.
    \item \textbf{Sign Flipping (SF)}: Sends $-2.5 \cdot g$, reversing the gradient
          direction and amplifying its magnitude.
    \item \textbf{FedPG Attack}: A sophisticated attack designed to pass the Byzantine
          filter while still degrading training.
    \item \textbf{Variance Attack (VA)}: Exploits gradient variance by injecting a
          gradient designed to increase the variance of the aggregated update.
    \item \textbf{Zero Gradient}: Sends a zero vector, simulating a dropped update
          without injecting noise.
    \item \textbf{Reward Flipping}: Negates all rewards during trajectory collection,
          causing the worker to learn the opposite of the intended behaviour.
\end{enumerate}

Byzantine workers are assigned by index: the first $B$ of $K$ workers are
Byzantine. This deterministic assignment ensures reproducibility and matches
the convention used in the FedPG-BR paper.

\section{Web Dashboard}

The web dashboard is a Flask~\cite{pallets2023flask} application with real-time
updates via Flask-SocketIO. It exposes a single experiment page at
\texttt{/experiment} where the user can:

\begin{itemize}
    \item Select environment, aggregation method, number of workers, number of
          Byzantine agents, attack type, and number of rounds.
    \item Expand an advanced settings panel to override hyperparameters
          ($\sigma$, $\delta$, learning rate, batch size, $\gamma$, etc.).
    \item Start and stop training with immediate effect.
    \item Monitor a live reward curve updated after every round.
    \item View per-round logs including the number of good agents retained by
          the Byzantine filter and the number of SCSG inner steps executed.
\end{itemize}

When running inside Docker, the dashboard communicates with a separately
spawned training container via the Docker API and a shared internal network.
When running locally, it launches \texttt{flower-simulation} as a subprocess
and streams its stdout through the WebSocket to the browser.

\section{Empirical Observations}

\subsection{CartPole-v1}

CartPole-v1 is the most reliable environment in this benchmark. Episodes are
capped at 500 steps regardless of policy quality, so per-round wall-clock time
remains bounded throughout training. The FedPG-BR Byzantine filter operates
stably because gradient statistics do not change dramatically as the policy
improves --- a solved CartPole policy still collects trajectories of similar
length. With $\sigma = 0.06$ and no Byzantine workers, FedPG-BR converges to
near-maximum reward within the first 50--100 rounds.

\subsection{LunarLander-v3 and the Filter Degradation Problem}

LunarLander-v3 reveals a fundamental tension between improving policy quality
and Byzantine filter stability. Early in training, the lander crashes quickly,
producing short episodes and low-variance gradients. The filter retains most
workers. As the policy improves, episodes become longer (the lander attempts
to land rather than crash immediately), producing higher-variance gradients.
The filter, parameterised with $\sigma = 0.07$, begins to classify honest
workers as Byzantine --- not because they are malicious, but because their
gradient norms now fall outside the tight threshold calibrated to early-training
statistics.

We observed the number of retained good agents drop from 10 to 5 within 20
rounds on LunarLander even in the absence of any Byzantine workers. This
over-filtering effect stalls learning: the aggregated gradient is computed from
an increasingly small subset of workers, reducing its quality and slowing
convergence.

This observation has a direct practical implication: the $\sigma$ parameter
must be tuned to the expected gradient variance at the intended operating point,
not at initialisation. Alternatively, SVRPG (which has no Byzantine filter)
avoids this problem entirely at the cost of Byzantine robustness.

\subsection{SCSG Ratio Check and Effective Variance Reduction}

A second observation concerns the SCSG inner loop. The importance ratio check
$0.995 \le \rho \le 1.005$ is extremely tight. In practice, after even a single
gradient step the policy parameters shift enough to push the importance ratio
outside this window, causing the loop to terminate with \texttt{scsg\_steps = 0}
or \texttt{1} in most rounds. The effective behaviour of SVRPG in this regime
is therefore indistinguishable from GOMDP --- plain gradient averaging with a
single update step.

This explains why LunarLander performance plateaued at approximately $+100$
reward with SVRPG: the variance reduction mechanism was not activating, and the
system was performing basic policy gradient with no additional signal.

\subsection{HalfCheetah-v5}

The HalfCheetah-v5 environment tests continuous control with a 17-dimensional
state space and 6-dimensional action space. MuJoCo episodes are time-limited
to 1000 steps regardless of performance, avoiding the episode-length growth
problem seen in LunarLander. Initial experiments show FedPG-BR starting at
approximately $-185$ reward and improving steadily over 208 rounds, with the
Byzantine filter behaving more stably due to the higher $\sigma = 0.9$ setting
appropriate for continuous control.
